% ===== PRÉAMBULE CANONIQUE — à coller VERBATIM en tête de CHAQUE .tex =====
\documentclass[11pt,a4paper]{article}
\usepackage[utf8]{inputenc}\usepackage[T1]{fontenc}\usepackage{lmodern}
\usepackage[french]{babel}
\usepackage{amsmath,amssymb,mathtools}\usepackage{icomma}
\usepackage{siunitx}\sisetup{locale=FR}  % \qty{}{}, \si{}, \ang{}
\usepackage{xcolor}\usepackage{colortbl}
\usepackage{tikz}\usepackage{pgfplots}\pgfplotsset{compat=1.18}
\usepackage[siunitx,europeanresistors]{circuitikz} % circuits (élec.)
\usepackage[most]{tcolorbox}\usepackage{enumitem}\usepackage{array,booktabs}
\usepackage[a4paper,margin=2.2cm]{geometry}
\usetikzlibrary{arrows.meta,calc,angles,quotes,positioning,patterns,%
  backgrounds,shapes.geometric,decorations.pathmorphing,decorations.markings,intersections}
\definecolor{primary}{RGB}{222,49,99}\definecolor{primarydark}{RGB}{150,22,55}
\definecolor{defcol}{RGB}{0,114,178}\definecolor{methcol}{RGB}{52,92,148}
\definecolor{excol}{RGB}{0,135,90}\definecolor{attncol}{RGB}{200,40,55}
\tcbset{boxbase/.style={enhanced,breakable,boxrule=0.9pt,arc=2.5pt,
  left=3mm,right=3mm,top=2.3mm,bottom=2.3mm,fonttitle=\bfseries,coltitle=white}}
% --- boîtes TUTORIEL (noms EXACTS, ne pas renommer) ---
\newtcolorbox{cle}[1][]{boxbase,colback=defcol!6,colframe=defcol,
  title={\textsc{À retenir}\ifx\relax#1\relax\relax\else\textmd{ : \itshape #1}\fi}}
\newtcolorbox{methode}[1][]{boxbase,colback=methcol!7,colframe=methcol,sharp corners,
  title={\textsc{Méthode}\ifx\relax#1\relax\relax\else\textmd{ --- \itshape #1}\fi}}
\newtcolorbox{exemplebox}[1][]{boxbase,colback=excol!5,colframe=excol,
  title={\textsc{Exemple}\ifx\relax#1\relax\relax\else\textmd{ : \itshape #1}\fi}}
\newtcolorbox{piege}[1][]{boxbase,colback=attncol!6,colframe=attncol,
  title={\textsc{Piège}\ifx\relax#1\relax\relax\else\textmd{ : \itshape #1}\fi}}
\newtcolorbox{remarque}[1][]{boxbase,colback=black!4,colframe=black!45,
  title={\textsc{Remarque}\ifx\relax#1\relax\relax\else\textmd{ : \itshape #1}\fi}}
% --- boîtes EXERCICE / CORRIGÉ (noms EXACTS, auto-comptées) ---
\newtcolorbox[auto counter]{exo}[1][]{boxbase,colback=primary!4,colframe=primary,
  title={\textsc{Exercice}~\thetcbcounter\ifx\relax#1\relax\relax\else\textmd{ --- \itshape #1}\fi}}
\newtcolorbox[auto counter]{corrige}[1][]{boxbase,colback=excol!4,colframe=excol,
  title={\textsc{Corrigé}~\thetcbcounter\ifx\relax#1\relax\relax\else\textmd{ --- \itshape #1}\fi}}
\newcommand{\vv}[1]{\vec{#1}}\newcommand{\ds}{\displaystyle}
\newcommand{\R}{\mathbb{R}}\newcommand{\N}{\mathbb{N}}\newcommand{\Z}{\mathbb{Z}}
% (un \usepackage{fancyhdr}+en-tête est possible ; il est ignoré côté web)
% =========================================================================
\begin{document}
\sisetup{per-mode=symbol}

\begin{corrige}[de l'électron à la charge]
On applique $q = N\,e$ ; perdre des électrons $\Rightarrow q>0$, en gagner $\Rightarrow q<0$.
\begin{enumerate}[label=\alph*)]
  \item Corps ayant \emph{perdu} des électrons : positif.
        \[ q = +N e = \num{5e12}\times\num{1.6e-19} = \SI{8e-7}{\coulomb} = \SI{0.8}{\micro\coulomb}. \]
  \item Corps ayant \emph{gagné} des électrons : négatif.
        \[ q = -N e = -\,\num{2.5e13}\times\num{1.6e-19} = \SI{-4e-6}{\coulomb} = \SI{-4}{\micro\coulomb}. \]
\end{enumerate}
\end{corrige}

\begin{corrige}[compter les électrons en excès]
\begin{enumerate}[label=\alph*)]
  \item \[ N = \frac{|q|}{e} = \frac{\num{8e-9}}{\num{1.6e-19}} = 5\times 10^{10}\ \text{électrons}. \]
  \item $q<0$ : la sphère porte un \textbf{excès} d'électrons ($\num{5e10}$ électrons de trop par
        rapport aux protons).
\end{enumerate}
\end{corrige}

\begin{corrige}[conducteur ou isolant ?]
\begin{enumerate}[label=\alph*)]
  \item \textbf{Conducteurs} : cuivre, argent, aluminium (métaux, charges libres). \quad
        \textbf{Isolants} : verre, plastique, ébonite (charges liées).
  \item Seuls les \textbf{isolants} (verre, plastique, ébonite) se chargent durablement par
        frottement : la charge y reste localisée. Un métal nu redistribue et écoule aussitôt toute
        charge, il ne se charge donc pas par frottement.
\end{enumerate}
\end{corrige}

\begin{corrige}[électrisation par contact]
\begin{enumerate}[label=\alph*)]
  \item Après contact, les deux sphères sont chargées \textbf{positivement}. Les électrons de la
        sphère neutre passent vers la sphère positive (déficitaire) jusqu'à égaliser : deux corps
        identiques se partagent la charge, donc chacun porte $+q/2$, de \emph{même signe}.
  \item Non : aucune charge n'est créée, on ne fait que \emph{répartir} les électrons existants.
        C'est le \textbf{principe de conservation de la charge}.
\end{enumerate}
\end{corrige}

\begin{corrige}[influence : le conducteur reste neutre]
\begin{enumerate}[label=\alph*)]
  \item Le bâton négatif \emph{repousse} les électrons libres de la barre : ils fuient vers
        l'extrémité \textbf{éloignée} (à droite), qui devient négative. L'extrémité \textbf{proche}
        (à gauche, face au bâton), appauvrie en électrons, devient positive.
  \item La barre reste \textbf{globalement neutre} : elle n'a ni gagné ni perdu d'électrons, ses
        charges se sont seulement \emph{déplacées} (polarisation). Sa charge totale est nulle.
\end{enumerate}
\end{corrige}

\end{document}
